\input{_preamble}
\begin{document}
%% Relative references with a sub-example part.  \Last[b] is linguex's
%% \complexExNo: the part belongs INSIDE the parentheses -- "(1b)", never
%% "(1)[b]" -- joined by \firstrefdash, so that a relative reference and a
%% \ref to a \sublabel spell the same example the same way.  Before the
%% commands took an optional argument at all, the bracket group simply fell
%% through into the text, which compiles perfectly and reads as a typo.
\ex.\label{rr:one}
\a.\sublabel{rr:a} First.
\b. Second.
\z.

BARE \Last{} SUB \Last[b] REF \ref{rr:a} PSUB \pLast[b]
NEXTSUB \Next[c] NNEXTSUB \NNext[d]

%% "\Last shows" must print "(1) shows", not "(1)shows": these are control
%% words, so TeX's tokenizer eats the following space long before any macro
%% could put it back, and linguex ends every one of these references in
%% \xspace for exactly that reason.  \xspace and not \space, because the
%% space must NOT appear before punctuation -- XSPCOMMA pins that half.
XSP \Last after PXSP \pLast after XSPARG \Last[b] after
XSPNEXT \Next after XSPCOMMA \Last, comma

\ex. Second example.

LLASTSUB \LLast[a] PNEXTSUB \pNext[e] PLLASTSUB \pLLast[f]

%% Inside a footnote the sub part rides on the footnote series (ii a),
%% while \TextNext keeps pointing at the main one -- the argument must not
%% disturb either branch.
FN\footnote{TEXTNEXTSUB \TextNext[a] PTEXTNEXTSUB \pTextNext[b]
\ex. Footnote one.

\ex. Footnote two.

FNSUB \Last[a] FNPSUB \pLast[a] FNLL \LLast[b] FNNEXT \Next[c]} done.

\ex. Target of TextNext.
\end{document}
